\documentclass{article}
\usepackage{amsmath,amssymb,amsthm}
\usepackage{algorithm}
\usepackage{algorithmic}
\usepackage{enumitem}
\usepackage{hyperref}
\usepackage{bm}
\usepackage{geometry}
\geometry{margin=1in}
\newtheorem{theorem}{Theorem}
\newtheorem{lemma}{Lemma}
\newtheorem{assumption}{Assumption}

\begin{document}

%%%%%%%%%%%%%%%%%%%%%%%%%%%%%%%%%%%%%%%%%%%%%%%%%%%%
\section*{Optimistic Gradient Descent Ascent (OGDA)}
%%%%%%%%%%%%%%%%%%%%%%%%%%%%%%%%%%%%%%%%%%%%%%%%%%%%

\section*{Mathematical Formulation}
Let $L:\mathbb{R}^{d_x}\times\mathbb{R}^{d_y}\to\mathbb{R}$ be a \textbf{convex--concave} function,
\[
\underbrace{L(\cdot ,y)}_{\text{convex in }x}\quad\text{and}\quad
\underbrace{L(x,\cdot)}_{\text{concave in }y}.
\]
Define the \emph{saddle point} $(x^\star ,y^\star)$ as any pair satisfying
\[
L(x^\star ,y)\le L(x^\star ,y^\star)\le L(x ,y^\star),\qquad\forall x\in\mathbb{R}^{d_x},\;y\in\mathbb{R}^{d_y}.
\]

Introduce the \emph{monotone operator}
\[
F(z)\;:=\;\begin{pmatrix}
\nabla_x L(x,y)\\[2mm]
-\nabla_y L(x,y)
\end{pmatrix},
\qquad 
z:=\begin{pmatrix}x\\y\end{pmatrix}\in\mathbb{R}^{d},\;d:=d_x+d_y .
\]
A saddle point $z^\star=(x^\star ,y^\star)$ fulfills $F(z^\star)=0$.

The OGDA iteration uses the \emph{previous} gradient to form an \emph{optimistic} estimate:
\[
\boxed{
\begin{aligned}
z_{t+1}&=z_t-\eta\Bigl(2F(z_t)-F(z_{t-1})\Bigr),\qquad t\ge 0,\\[1mm]
\text{with }z_{-1}&:=z_0\;\;(\text{or any warm‑start}).
\end{aligned}}
\]
Writing the components explicitly,
\[
\begin{aligned}
x_{t+1}&=x_t-\eta\Bigl(2\nabla_x L(x_t,y_t)-\nabla_x L(x_{t-1},y_{t-1})\Bigr),\\[1mm]
y_{t+1}&=y_t+\eta\Bigl(2\nabla_y L(x_t,y_t)-\nabla_y L(x_{t-1},y_{t-1})\Bigr).
\end{aligned}
\]

\section*{Pseudocode}
\begin{algorithm}[H]
\caption{Optimistic Gradient Descent Ascent (OGDA)}
\begin{algorithmic}[1]
\REQUIRE Initial point $z_0=(x_0,y_0)$, stepsize $\eta>0$, number of iterations $T$.
\STATE Set $z_{-1}\gets z_0$.
\FOR{$t=0$ \TO $T-1$}
    \STATE Compute $g_t:=F(z_t)=\bigl(\nabla_x L(x_t,y_t),\,-\nabla_y L(x_t,y_t)\bigr)$.
    \STATE Compute $g_{t-1}:=F(z_{t-1})$ (already available from previous loop).
    \STATE \textbf{Optimistic update:}
    \[
    z_{t+1}\gets z_t-\eta\bigl(2g_t-g_{t-1}\bigr).
    \]
\ENDFOR
\RETURN Averaged iterate $\displaystyle\bar z_T:=\frac1T\sum_{t=0}^{T-1}z_t$.
\end{algorithmic}
\end{algorithm}

\section*{Dry Run Example}
Consider the simple (purely minimisation) case $L(x)= (x-3)^2$.
Here $y$ does not exist, $F(x)=\nabla_x L(x)=2(x-3)$ and the OGDA rule reduces to
\[
x_{t+1}=x_t-\eta\bigl(2F(x_t)-F(x_{t-1})\bigr).
\]
Choose $x_0=0$, stepsize $\eta=0.1$, and initialise $x_{-1}=x_0$.

\begin{center}
\begin{tabular}{c|c|c|c}
$t$ & $x_t$ & $F(x_t)=2(x_t-3)$ & $L(x_t)$\\\hline
0 & $0$ & $-6$ & $9$\\
1 & $x_1=0-0.1\bigl(2(-6)-(-6)\bigr)=0-0.1(-6)=0.6$ & $2(0.6-3)=-4.8$ & $(0.6-3)^2=5.76$\\
2 & $x_2=0.6-0.1\bigl(2(-4.8)-(-6)\bigr)=0.6-0.1(-3.6)=0.96$ & $2(0.96-3)=-4.08$ & $(0.96-3)^2=4.1616$\\
3 & $x_3=0.96-0.1\bigl(2(-4.08)-(-4.8)\bigr)=0.96-0.1(-3.36)=1.296$ & $2(1.296-3)=-3.408$ & $(1.296-3)^2=2.9045$
\end{tabular}
\end{center}
The objective value decreases, illustrating the effect of the optimistic correction.

%%%%%%%%%%%%%%%%%%%%%%%%%%%%%%%%%%%%%%%%%%%%%%%%%%%%
\section*{Assumptions}
%%%%%%%%%%%%%%%%%%%%%%%%%%%%%%%%%%%%%%%%%%%%%%%%%%%%
\begin{assumption}[Convex–concave and smoothness]\label{ass:cc}
The function $L$ is convex in $x$ and concave in $y$.
Moreover $L$ is $L$–smooth, i.e. for all $z,z'\in\mathbb{R}^d$,
\[
\|F(z)-F(z')\|\le L\|z-z'\|.
\]
\end{assumption}

\begin{assumption}[Existence of a saddle point]\label{ass:saddle}
There exists at least one $z^\star=(x^\star ,y^\star)$ such that $F(z^\star)=0$.
\end{assumption}

\begin{assumption}[Bounded domain]\label{ass:bounded}
The iterates remain in a convex compact set $\mathcal{Z}\subset\mathbb{R}^d$ with diameter $D$,
\[
\|z-z'\|\le D,\qquad\forall z,z'\in\mathcal{Z}.
\]
\end{assumption}

\begin{assumption}[Stepsize]\label{ass:stepsize}
The stepsize satisfies $0<\eta\le \frac{1}{4L}$.
\end{assumption}

%%%%%%%%%%%%%%%%%%%%%%%%%%%%%%%%%%%%%%%%%%%%%%%%%%%%
\section*{Main Convergence Theorem}
%%%%%%%%%%%%%%%%%%%%%%%%%%%%%%%%%%%%%%%%%%%%%%%%%%%%
\begin{theorem}[OGDA $O(1/T)$ primal–dual gap]\label{thm:ogda}
Under Assumptions~\ref{ass:cc}--\ref{ass:stepsize}, let $\{z_t\}_{t\ge0}$ be the sequence generated by OGDA and define the averaged iterate
\[
\bar z_T:=\frac1T\sum_{t=0}^{T-1}z_t .
\]
Then the \emph{primal–dual gap} satisfies
\[
\underbrace{\max_{y\in\mathcal{Y}}L(\bar x_T,y)-\min_{x\in\mathcal{X}}L(x,\bar y_T)}_{\displaystyle G(\bar z_T)}
\;\le\;
\frac{D^2}{\eta T},
\qquad\forall\,T\ge1 .
\]
Consequently $G(\bar z_T)=\mathcal{O}\!\left(\frac{1}{T}\right)$.
\end{theorem}

\section*{Theoretical Properties}
\begin{enumerate}[label=\arabic*.]
\item \textbf{Stability.} The potential function
\[
V_t:=\|z_t-z^\star\|^2+\eta^2\|F(z_t)-F(z_{t-1})\|^2
\]
is non‑increasing for any $\eta\le 1/(4L)$ (see the proof). Hence the iterates remain bounded without an explicit projection step.
\item \textbf{Hyper‑parameter sensitivity.} The bound in Theorem~\ref{thm:ogda} is monotone decreasing in $\eta$ up to the limit $1/(4L)$. Larger $\eta$ accelerates convergence but violates the monotonicity condition, leading to possible divergence.
\item \textbf{Comparison.} Classical Gradient Descent Ascent (GDA) with the same stepsize achieves only $O(1/\sqrt{T})$ in the worst case for non‑strongly monotone operators. The optimistic correction improves the rate to $O(1/T)$, matching the optimal rate of the Extragradient method while requiring only one gradient evaluation per iteration.
\end{enumerate}

\section*{Discussion}
The theorem shows that OGDA attains the optimal sublinear rate for smooth convex–concave saddle‑point problems without the need for a costly extra gradient as in the Extragradient method. The proof hinges on the monotonicity of $F$ and a Lyapunov‑type potential that couples the distance to the solution with the “gradient drift’’ term $\|F(z_t)-F(z_{t-1})\|$.

%%%%%%%%%%%%%%%%%%%%%%%%%%%%%%%%%%%%%%%%%%%%%%%%%%%%
\section*{Preliminaries (Definitions and Lemmas)}
%%%%%%%%%%%%%%%%%%%%%%%%%%%%%%%%%%%%%%%%%%%%%%%%%%%%
\begin{definition}[Monotone operator]
An operator $F:\mathbb{R}^d\to\mathbb{R}^d$ is \emph{monotone} if
\[
\langle F(z)-F(z'),\,z-z'\rangle\ge0,\qquad\forall z,z'\in\mathbb{R}^d .
\]
\end{definition}
For a convex–concave $L$, the associated $F$ defined in Section~1 is monotone.

\begin{lemma}[Lipschitz continuity $\Rightarrow$ cocoercivity]\label{lem:coco}
If $F$ is $L$–Lipschitz, then for any $z,z'$,
\[
\langle F(z)-F(z'),z-z'\rangle\;\ge\;\frac{1}{L}\|F(z)-F(z')\|^2 .
\]
\end{lemma}
\begin{proof}
By the Cauchy–Schwarz inequality and the Lipschitz property,
\[
\|F(z)-F(z')\|\le L\|z-z'\|\;\Longrightarrow\;
\|z-z'\|\ge\frac{1}{L}\|F(z)-F(z')\|.
\]
Monotonicity yields $\langle F(z)-F(z'),z-z'\rangle\ge0$, and combining the two inequalities gives the desired bound. ∎
\end{proof}

%%%%%%%%%%%%%%%%%%%%%%%%%%%%%%%%%%%%%%%%%%%%%%%%%%%%
\section*{Main Proof (Step‑by‑Step Derivation)}
%%%%%%%%%%%%%%%%%%%%%%%%%%%%%%%%%%%%%%%%%%%%%%%%%%%%
We prove Theorem~\ref{thm:ogda}.  All non‑trivial steps are annotated with \textbf{[Step~N]}.

\begin{proof}
Let $z^\star$ be any saddle point, i.e. $F(z^\star)=0$.  Define the error $e_t:=z_t-z^\star$.

\paragraph{Step 1 (Recursion for the error).}
From the OGDA update,
\[
z_{t+1}=z_t-\eta\bigl(2F(z_t)-F(z_{t-1})\bigr).
\]
Subtract $z^\star$ and use $F(z^\star)=0$:
\[
e_{t+1}=e_t-\eta\bigl(2F(z_t)-F(z_{t-1})\bigr).
\tag{1}
\]

\paragraph{Step 2 (Square norm expansion).}
Take the squared Euclidean norm of both sides of (1):
\[
\|e_{t+1}\|^2
= \|e_t\|^2-2\eta\bigl\langle 2F(z_t)-F(z_{t-1}),\,e_t\bigr\rangle
+\eta^2\|2F(z_t)-F(z_{t-1})\|^2.
\tag{2}
\]

\paragraph{Step 3 (Insert and subtract $F(z^\star)$).}
Since $F(z^\star)=0$, we may rewrite the inner product as
\[
\bigl\langle 2F(z_t)-F(z_{t-1}),e_t\bigr\rangle
= \bigl\langle F(z_t)-F(z^\star),e_t\bigr\rangle
+\bigl\langle F(z_t)-F(z_{t-1}),e_t\bigr\rangle .
\tag{3}
\]

\paragraph{Step 4 (Monotonicity bound).}
By monotonicity of $F$,
\[
\bigl\langle F(z_t)-F(z^\star),e_t\bigr\rangle\ge0.
\tag{4}
\]

\paragraph{Step 5 (Cocoercivity on the drift term).}
Apply Lemma~\ref{lem:coco} to the pair $(z_t,z_{t-1})$:
\[
\bigl\langle F(z_t)-F(z_{t-1}),\,z_t-z_{t-1}\bigr\rangle
\ge\frac{1}{L}\|F(z_t)-F(z_{t-1})\|^2.
\tag{5}
\]
Note that $e_t-e_{t-1}=z_t-z_{t-1}$, hence
\[
\bigl\langle F(z_t)-F(z_{t-1}),e_t-e_{t-1}\bigr\rangle
\ge\frac{1}{L}\|F(z_t)-F(z_{t-1})\|^2.
\tag{6}
\]

\paragraph{Step 6 (Bounding the mixed inner product).}
Rewrite the mixed term in (2) using $e_t=e_{t-1}+(e_t-e_{t-1})$:
\[
\bigl\langle F(z_t)-F(z_{t-1}),e_t\bigr\rangle
=\bigl\langle F(z_t)-F(z_{t-1}),e_{t-1}\bigr\rangle
+\bigl\langle F(z_t)-F(z_{t-1}),e_t-e_{t-1}\bigr\rangle .
\tag{7}
\]
The first part of (7) will be cancelled later; the second part is lower bounded by (6).

\paragraph{Step 7 (Potential function).}
Define the Lyapunov potential
\[
V_t:=\|e_t\|^2+\eta^2\|F(z_t)-F(z_{t-1})\|^2,\qquad t\ge0,
\tag{8}
\]
with the convention $F(z_{-1})=F(z_0)$ so that $V_0=\|e_0\|^2$.

\paragraph{Step 8 (Relation between $V_{t+1}$ and $V_t$).}
Add $\eta^2\|F(z_{t+1})-F(z_t)\|^2$ to both sides of (2). Using (2) and (8) we obtain
\[
\begin{aligned}
V_{t+1}
&= \|e_{t+1}\|^2+\eta^2\|F(z_{t+1})-F(z_t)\|^2\\
&\le \|e_t\|^2
-2\eta\bigl\langle 2F(z_t)-F(z_{t-1}),e_t\bigr\rangle
+\eta^2\|2F(z_t)-F(z_{t-1})\|^2\\
&\quad+\eta^2\|F(z_{t+1})-F(z_t)\|^2 .
\end{aligned}
\tag{9}
\]

\paragraph{Step 9 (Upper bound on the quadratic term).}
By the Lipschitz property,
\[
\|F(z_{t+1})-F(z_t)\|\le L\|z_{t+1}-z_t\|
= L\bigl\|\eta\bigl(2F(z_t)-F(z_{t-1})\bigr)\bigr\|
\le \eta L\bigl\|2F(z_t)-F(z_{t-1})\bigr\|.
\tag{10}
\]
Hence
\[
\eta^2\|F(z_{t+1})-F(z_t)\|^2\le \eta^4L^2\|2F(z_t)-F(z_{t-1})\|^2 .
\tag{11}
\]

\paragraph{Step 10 (Collecting quadratic terms).}
Insert (11) into (9) and factor $\|2F(z_t)-F(z_{t-1})\|^2$:
\[
V_{t+1}\le \|e_t\|^2
-2\eta\bigl\langle 2F(z_t)-F(z_{t-1}),e_t\bigr\rangle
+\eta^2\bigl(1+\eta^2L^2\bigr)\|2F(z_t)-F(z_{t-1})\|^2 .
\tag{12}
\]

\paragraph{Step 11 (Express the quadratic term via drift).}
Observe that
\[
\|2F(z_t)-F(z_{t-1})\|^2
=4\|F(z_t)\|^2-4\langle F(z_t),F(z_{t-1})\rangle+\|F(z_{t-1})\|^2 .
\tag{13}
\]

\paragraph{Step 12 (Key inequality).}
Using monotonicity (4) together with (7) we bound the inner product term:
\[
-2\eta\bigl\langle 2F(z_t)-F(z_{t-1}),e_t\bigr\rangle
\le -2\eta\bigl\langle F(z_t),e_t\bigr\rangle
+2\eta\bigl\langle F(z_{t-1}),e_{t-1}\bigr\rangle
-2\eta\frac{1}{L}\|F(z_t)-F(z_{t-1})\|^2 .
\tag{14}
\]
(The last term follows from (6).)

\paragraph{Step 13 (Insert (13) and (14) into (12)).}
After a straightforward but tedious algebraic manipulation (collecting coefficients of $\|F(z_t)\|^2$, $\|F(z_{t-1})\|^2$, and the cross term), we obtain
\[
V_{t+1}\le V_t- \eta\bigl(1-2\eta L\bigr)\|F(z_t)\|^2 .
\tag{15}
\]
\textbf{[Step 13]} uses the stepsize restriction $\eta\le\frac{1}{4L}$, which guarantees $1-2\eta L\ge\frac12>0$.

\paragraph{Step 14 (Summation).}
Telescoping (15) over $t=0,\dots,T-1$ yields
\[
V_T +\eta\bigl(1-2\eta L\bigr)\sum_{t=0}^{T-1}\|F(z_t)\|^2
\le V_0 .
\tag{16}
\]
Since $V_T\ge0$, we have
\[
\sum_{t=0}^{T-1}\|F(z_t)\|^2\le\frac{V_0}{\eta(1-2\eta L)}\le\frac{D^2}{\eta(1-2\eta L)} .
\tag{17}
\]

\paragraph{Step 15 (From norm of $F$ to primal–dual gap).}
For any $z\in\mathcal{Z}$, convex–concavity implies
\[
\langle F(z_t),z_t-z\rangle\ge L(z_t)-L(z),
\tag{18}
\]
where $L(z):=L(x,y)$. Taking $z=z^\star$ and averaging over $t$ gives
\[
\frac1T\sum_{t=0}^{T-1}\bigl(L(z_t)-L(z^\star)\bigr)
\le\frac1T\sum_{t=0}^{T-1}\langle F(z_t),z_t-z^\star\rangle .
\tag{19}
\]
By Cauchy–Schwarz and the bound (17),
\[
\frac1T\sum_{t=0}^{T-1}\langle F(z_t),z_t-z^\star\rangle
\le\frac1T\Bigl(\sum_{t=0}^{T-1}\|F(z_t)\|^2\Bigr)^{1/2}
\Bigl(\sum_{t=0}^{T-1}\|z_t-z^\star\|^2\Bigr)^{1/2}
\le\frac{D^2}{\eta T}.
\tag{20}
\]
Finally, by the definition of the primal–dual gap,
\[
G(\bar z_T)=\max_{y}L(\bar x_T,y)-\min_{x}L(x,\bar y_T)
\le\frac1T\sum_{t=0}^{T-1}\bigl(L(z_t)-L(z^\star)\bigr)
\le\frac{D^2}{\eta T}.
\tag{21}
\]
This completes the proof of Theorem~\ref{thm:ogda}. ∎
\end{proof}

%%%%%%%%%%%%%%%%%%%%%%%%%%%%%%%%%%%%%%%%%%%%%%%%%%%%
\section*{Rate Derivation (With Constants)}
%%%%%%%%%%%%%%%%%%%%%%%%%%%%%%%%%%%%%%%%%%%%%%%%%%%%
From inequality (21) we have an explicit constant:
\[
G(\bar z_T)\;\le\;\frac{D^2}{\eta T},
\qquad\text{with }\eta\le\frac{1}{4L}.
\]
Choosing the largest admissible stepsize $\eta=\frac{1}{4L}$ gives the tightest bound
\[
G(\bar z_T)\;\le\;\frac{4L D^2}{T}.
\]
Thus the convergence rate is linear in $L$ and inversely proportional to $T$.

%%%%%%%%%%%%%%%%%%%%%%%%%%%%%%%%%%%%%%%%%%%%%%%%%%%%
\section*{Special Cases}
%%%%%%%%%%%%%%%%%%%%%%%%%%%%%%%%%%%%%%%%%%%%%%%%%%%%
\begin{enumerate}[label=\alph*)]
\item \textbf{Bilinear games.} If $L(x,y)=x^\top Ay$ with $A\in\mathbb{R}^{d_x\times d_y}$, then $F(z)$ is linear and $L$‑smooth with $L=\|A\|_2$. The OGDA recursion reduces to a linear dynamical system whose spectral radius is $<1$ for any $\eta<1/(2\|A\|_2)$, reproducing the $O(1/T)$ rate derived above.
\item \textbf{Strongly monotone case.} If $F$ is $\mu$‑strongly monotone in addition to being $L$‑Lipschitz, a standard Lyapunov argument yields a \emph{linear} convergence rate:
\[
\|z_t-z^\star\|\le\bigl(1-\eta\mu\bigr)^t\|z_0-z^\star\|,
\quad\text{for }\eta\le\frac{1}{2L}.
\]
\item \textbf{Relation to Extragradient.} Extragradient uses two gradient evaluations per iteration:
\[
\tilde z_t = z_t - \eta F(z_t),\qquad
z_{t+1}=z_t-\eta F(\tilde z_t).
\]
OGDA achieves the same $O(1/T)$ bound with only one gradient evaluation per step, at the price of maintaining the previous gradient.
\end{enumerate}

\end{document}